% songs.tex — Song Annotation Schema in LaTeX
% Superset format covering: lyrics, chords, brackets, annotations,
% struck text, side cues, highlights, circled refs, structure markers.
%
% Compile with: xelatex songs.tex  (for fontspec)
% Or:           pdflatex songs.tex (remove fontspec block, use lmodern)

\documentclass[11pt, a4paper]{article}

% ── Packages ──────────────────────────────────────────────────────────
\usepackage{fontspec}           % XeLaTeX: system fonts
\setmainfont{EB Garamond}       % or: Georgia, Palatino — fallback below
\setmonofont{Inconsolata}       % for chords, annotations
% If compiling with pdflatex, replace the three lines above with:
%   \usepackage[T1]{fontenc}
%   \usepackage{lmodern}

\usepackage[margin=2.5cm, right=4.5cm]{geometry}  % wide right margin for side cues
\usepackage{xcolor}
\usepackage{soul}               % \st{} strikethrough, \hl{} highlight
\usepackage{ulem}               % \sout{} — better strikethrough
\usepackage{marginnote}         % \marginnote{} for side cues
\usepackage{mdframed}           % bracket-group box and side-box
\usepackage{enumitem}
\usepackage{tikz}               % circled references
\usepackage{array, booktabs}    % cue-sheet tables
\usepackage{ragged2e}
\usepackage{microtype}
\usepackage{hyperref}
\hypersetup{colorlinks=true, linkcolor=black}

\normalem  % don't let ulem hijack \emph

% ── Colour palette ───────────────────────────────────────────────────
\definecolor{chordblue}   {HTML}{1a6080}
\definecolor{strikegray}  {HTML}{9a8060}
\definecolor{annotred}    {HTML}{b04000}
\definecolor{blueinkcol}  {HTML}{1a4080}
\definecolor{sidegreen}   {HTML}{5a6030}
\definecolor{stagebrn}    {HTML}{5a4020}
\definecolor{bracketbrn}  {HTML}{5a4020}
\definecolor{altpurple}   {HTML}{5a4060}

% Highlights
\definecolor{hlgreen}  {HTML}{b8f070}
\definecolor{hlpink}   {HTML}{ff8fb0}
\definecolor{hlblue}   {HTML}{70d4f8}
\definecolor{hlorange} {HTML}{ffc060}
\definecolor{hlyellow} {HTML}{f0f050}

% ── Schema commands ───────────────────────────────────────────────────

% Chord symbol (appears inline before the syllable it falls on)
\newcommand{\chord}[1]{%
  {\ttfamily\small\color{chordblue}\textsuperscript{#1}}%
}

% Chord line (a full line of chords above a lyric line)
\newcommand{\chordline}[1]{%
  \par\noindent{\ttfamily\small\color{chordblue}#1}\par\vspace{-0.6\baselineskip}%
}

% Struck-through text (original, crossed out)
\newcommand{\struck}[1]{{\color{strikegray}\sout{#1}}}

% Inserted replacement text (follows a \struck)
\newcommand{\ins}[1]{{\bfseries #1}}

% Highlight macros
\newcommand{\hlgreen}[1]{\colorbox{hlgreen}{\small\ttfamily\color{black!80}#1}}
\newcommand{\hlpink}[1]{\colorbox{hlpink}{\small\ttfamily\color{black!80}#1}}
\newcommand{\hlblue}[1]{\colorbox{hlblue}{\small\ttfamily\color{black!80}#1}}
\newcommand{\hlorange}[1]{\colorbox{hlorange}{\small\ttfamily\color{black!80}#1}}
\newcommand{\hlyellow}[1]{\colorbox{hlyellow}{\small\ttfamily\color{black!80}#1}}

% Circled reference
\newcommand{\circref}[1]{%
  \tikz[baseline=(char.base)]{%
    \node[shape=circle, draw, inner sep=1.5pt, font=\ttfamily\scriptsize] (char) {#1};%
  }%
}

% Side cue (right margin, performance direction)
\newcommand{\sidecue}[1]{%
  \marginnote{\itshape\small\color{sidegreen}#1}%
}

% Stage direction (spoken aside, inline)
\newcommand{\stage}[1]{{\itshape\color{stagebrn}\small(#1)}}

% Blue-ink annotation (different hand/pass)
\newcommand{\blueink}[1]{%
  \par\noindent\hspace{1em}{\ttfamily\small\color{blueinkcol}\itshape ✎\enspace #1}\par%
}

% Cut line (marked ✕)
\newcommand{\cutline}[1]{%
  {\color{strikegray}\sout{#1}\enspace\texttt{\scriptsize[✕]}}%
}

% Arrow annotation
\newcommand{\arrowann}[1]{%
  \par\noindent\hspace{1em}{\color{annotred}\small → #1}\par%
}

% Marginal/corner note
\newcommand{\cornernote}[1]{%
  \marginnote{\ttfamily\scriptsize\color{gray}\itshape #1}%
}

% ── Block environments ────────────────────────────────────────────────

% verse: a verse or section block
\newenvironment{verse}%
  {\vspace{0.6em}\noindent}%
  {\vspace{0.6em}}

% bracket-group: lines grouped with a left bracket
\newmdenv[
  leftline=true, rightline=false, topline=false, bottomline=false,
  linecolor=bracketbrn, linewidth=2.5pt,
  innerleftmargin=10pt, innerrightmargin=0pt,
  innertopmargin=2pt, innerbottommargin=2pt,
  backgroundcolor=white
]{bracketgroup}

% annotation: handwritten marginal note
\newenvironment{annotation}%
  {\par\noindent\hspace{1em}\begin{minipage}{0.85\linewidth}%
   \ttfamily\small\color{annotred}\itshape ↳\enspace\ignorespaces}%
  {\end{minipage}\par}

% alt: alternative version of a line
\newenvironment{alt}%
  {\par\vspace{2pt}\noindent\hspace{1.2em}%
   \begin{minipage}{0.82\linewidth}%
   \small\itshape\color{altpurple}%
   \setlength{\leftskip}{6pt}%
   \setlength{\parindent}{0pt}}%
  {\end{minipage}\par\vspace{2pt}}

% side-box: boxed reference panel
\newmdenv[
  linecolor=strikegray, linewidth=1.5pt,
  backgroundcolor=yellow!5,
  innerleftmargin=10pt, innerrightmargin=8pt,
  innertopmargin=8pt, innerbottommargin=8pt,
  skipabove=6pt, skipbelow=6pt
]{sidebox}

% structure-label: Intro, Outro, BUT, [A], etc.
\newcommand{\structlabel}[1]{%
  \vspace{0.8em}\noindent{\ttfamily\small\bfseries\MakeUppercase{#1}}\vspace{0.3em}\par%
}

% ── Title / section formatting ─────────────────────────────────────────
\usepackage{titlesec}
\titleformat{\section}{\large\bfseries}{}{0em}{}[\vspace{-0.3em}\rule{\textwidth}{0.8pt}]
\titleformat{\subsection}{\normalsize\itshape}{}{0em}{}

% ═══════════════════════════════════════════════════════════════════════
\begin{document}
% ═══════════════════════════════════════════════════════════════════════

\begin{center}
  {\LARGE\bfseries Song Transcriptions}\\[0.4em]
  {\small\itshape Trad. — with arrangement and annotation notes}
\end{center}

\tableofcontents
\newpage

% ──────────────────────────────────────────────────────────────────────
\section{Song 1 — (Reluctant Fisherman)}
\subsection*{Trad.}

\structlabel{[A]}

\begin{verse}
\begin{bracketgroup}
Oh me dad was a fisherman bold\\
And he lived til he grew old. \sidecue{For\ldots}\\
\blueink{For he'd open the panes \& pop out the flame / Just to see how the wind do blow}
And he'd often say to me\\
You'd be wise before you go. \sidecue{Do\ldots}
\end{bracketgroup}
\end{verse}

\hrulefill

\begin{verse}
\begin{bracketgroup}
When the North wind rough do blow\\
Then I lay right smug below. \sidecue{But\ldots}
\end{bracketgroup}
\end{verse}

\hrulefill

\begin{verse}
\begin{bracketgroup}
When the wind comes out of the East\\
It's no use to man nor beast. \sidecue{But\ldots}
\begin{alt}
East – wind out of the East\\
\struck{Poor} \ins{Yosh} wife – We'll regret if
\end{alt}
\end{bracketgroup}
\end{verse}

\hrulefill

\begin{verse}
\begin{bracketgroup}
When the wind's back in to the west\\
That'll come in rough at best. \sidecue{But\ldots}
\begin{alt}
\struck{North-wind \& so} West wind\\
\stage{A lady says /} \struck{poor wife} --- You'll regret if\\
\arrowann{\textbf{Fine wife}}
\end{alt}
\end{bracketgroup}
\end{verse}

\hrulefill

\begin{verse}
\begin{bracketgroup}
When the south wind soft do blow\\
Then there's not enough wind to go. \sidecue{Bu[t]\ldots}
\end{bracketgroup}
\end{verse}

\hrulefill

\begin{verse}
\begin{bracketgroup}
\cutline{North Wind}\\
\cutline{\struck{Royal Lady (not wife)}}\\
\struck{Poor wife} – You will pay
\end{bracketgroup}
\end{verse}

\hrulefill

\begin{verse}
\begin{bracketgroup}
And me poor wife says to me\\
We shall starve if you don't go. \sidecue{So\ldots}
\end{bracketgroup}
\begin{annotation}
\stage{Are you mad – knew wouldn't get out of it} — spoken/aside
\end{annotation}
\end{verse}

\hrulefill

\begin{verse}
Now all you fishermen bold\\
If you live till you grow old. \sidecue{Do\ldots}
\end{verse}


% ──────────────────────────────────────────────────────────────────────
\newpage
\section{Song 2 — The Two Sisters (Oh the Wind \& Rain)}
\subsection*{TRAD. — Peggy Seeger / Bay Whitaker / Fates}

\structlabel{Intro – our version}

\begin{verse}
\begin{bracketgroup}
\chordline{A \hspace{1.8em} B \hspace{2em} D}
Twas early one morning in the month of May\\[2pt]
\chordline{Em \hspace{3.8em} A}
\textit{Oh the wind and rain;}\\[2pt]
\chordline{A \hspace{4em} D}
Two sisters went fishing on a hot summer's day,\\[2pt]
\chordline{Em \hspace{1.8em} D \hspace{1.8em} A}
\textit{\stage{CRYING} Oh, the dreadful wind and rain.}
\end{bracketgroup}
\end{verse}

\begin{annotation}
\circref{A} section repeat label
\end{annotation}

\hrulefill

\begin{verse}
\begin{bracketgroup}
\chordline{A \hspace{4em} D \hspace{3em} Em -- A}
Two sweet sisters side by side / \textit{Oh\ldots}\\[2pt]
\chordline{A \hspace{5em} D \hspace{3.5em} Em -- D -- A}
Both of them want to be Johnny's bride / \textit{Oh\ldots} \stage{CRYING}
\end{bracketgroup}
\begin{alt}
a \struck{Good} \ins{Good} man's\\
or: the same man's
\end{alt}
\end{verse}

\hrulefill

\begin{verse}
\begin{bracketgroup}
\cutline{Johnny tells older sister it's not her}\\[2pt]
→ Johnny gave the young one a fine gold ring / \textit{Oh\ldots /}
\begin{annotation}
He then\ldots
\end{annotation}
Didn't give the other one anything / \textit{Oh\ldots}
\end{bracketgroup}
\end{verse}

\hrulefill

\begin{verse}
\begin{bracketgroup}
Two sweet sisters walking by the stream / \textit{Oh\ldots /}\\
One came behind \& pushed the other one in / \textit{Oh\ldots}
\end{bracketgroup}
\end{verse}

\hrulefill

\begin{verse}
\begin{bracketgroup}
She pushed her in to the river to drown / \textit{Oh\ldots /}\\
And watched her as she floated down / \textit{Oh\ldots}
\end{bracketgroup}
\end{verse}

\hrulefill

\begin{verse}
\begin{bracketgroup}
She floated on down to the miller's pond. / \textit{Oh\ldots /}\\
``Father, father, there comes a swan!'' / \textit{Oh\ldots}
\begin{annotation}
straight to → ``But it wasn't a swan'' — then straight into next verse
\end{annotation}
\end{bracketgroup}
\end{verse}

\hrulefill

\begin{verse}
\begin{bracketgroup}
The miller ran for his drifting hook / \textit{Oh\ldots /}\\
And brought that poor girl from the brook / \textit{Oh\ldots}
\end{bracketgroup}
\end{verse}

\hrulefill

\begin{verse}
\begin{bracketgroup}
\cutline{He laid her on the bank so dry / \textit{Oh\ldots /}}\\
\cutline{A fiddler man came walking by / \textit{Oh\ldots}}
\end{bracketgroup}
\begin{annotation}
both lines marked ✕ (possibly cut)
\end{annotation}
\end{verse}

\hrulefill

\begin{verse}
\begin{bracketgroup}
He \struck{saw that poor girl lying} \ins{found me / buried} there / \textit{Oh\ldots /}
\begin{annotation}
alt: ``found me'' / ``buried'' written above ``saw that'' / ``lying''
\end{annotation}
\cutline{He took thirty strands of her long yellow hair / \textit{Oh\ldots}}
\begin{annotation}
Fashioned from\ldots\ And that's exactly what he did
\end{annotation}
\end{bracketgroup}
\end{verse}

\hrulefill

\begin{verse}
\begin{bracketgroup}
He made a fiddle bow of her long yellow hair / \textit{Oh\ldots /}\\
He made fiddle pegs of her little finger bones / \textit{Oh\ldots}
\end{bracketgroup}
\end{verse}

\hrulefill

\begin{verse}
\begin{bracketgroup}
He made a fiddle of her little breast-bone / \textit{Oh\ldots /}\\
With a sound that could melt a heart of stone / \textit{Oh\ldots}
\begin{annotation}
He used to play other fiddles — it would only play\ldots
\end{annotation}
\end{bracketgroup}
\end{verse}

\hrulefill

\begin{verse}
And the only tune that fiddle would play / \textit{Oh\ldots}\\
The only tune that fiddle would play, was / \textit{Oh\ldots}
\begin{annotation}
song they'd sung together – older sister's face
\end{annotation}
\end{verse}

\medskip
\textit{Oh, the wind and rain.}

\textit{The only tune that the fiddle would play, was oh the dreadful wind and rain.}


% ──────────────────────────────────────────────────────────────────────
\newpage
\section{Song 3 — MacPherson's Lament}
\subsection*{Bk6 No.10}

\structlabel{Cue Sheet / Arrangement Notes}

\begin{tabular}{@{} p{6.5cm} p{4.5cm} l @{}}
\toprule
\textbf{Beat / Story Point} & \textbf{Tune Cue} & \textbf{Ref} \\
\midrule
Opening: Poem, This is his tale / tune
  & \hlgreen{3 McP's}
  & \circref{p1} \\
Boy being brought up — gypsy + accomplished musician / fiddler
  & \hlpink{\struck{8 Flowers}}
  & \\
Father dies – kicked out
  & MacPh's
  & \circref{p3} \\
Plays for band of gypsies
  & \hlpink{18 My Love}
  & \\
Capture of Jannie
  & \hlblue{18 part 2}
  & \\
Escape
  & \hlblue{Dashing W.S.}
  & \\
No escape
  & ---
  & \\
Last request
  & \hlyellow{McPs L}
  & \circref{3} \circref{P5} \\
Known as M's Lament = Play in key of\ldots
  & ---
  & \\
Outro → \textit{(time: 1:50)}
  & \hlorange{G with double stopping}
  & \\
\bottomrule
\end{tabular}

\begin{annotation}
Paul Anderson – Spotify
\end{annotation}

\begin{sidebox}
\textbf{\circref{8} Flowers of Edinburgh} — Last 2 lines\\[4pt]
\textbf{\circref{18}} \textit{(My Love She's But A Lassie Yet)}\\
\arrowann{bar 1–34}
\arrowann{bar 99–end}
Dashing White Sergeant\\
\textit{(partially cut off in source)}
\end{sidebox}


% ──────────────────────────────────────────────────────────────────────
\newpage
\section{Song 4 — Raggle Taggle}
\cornernote{top-right corner: ``muttering = +music''}

\structlabel{Cue Sheet / Arrangement Notes}

\begin{tabular}{@{} p{7.5cm} p{5cm} @{}}
\toprule
\textbf{Beat / Story Point} & \textbf{Tune Cue} \\
\midrule
Intro: Play \struck{together}
  & \hlyellow{RTG.} \\
Sense happiness
  & \hlorange{Coconut dance (a)} \\
Leaving her for a year
  & \hlpink{Hidden Valley} \\
\hspace{1em} ↳ Letter – killed
  & \\
\hspace{1em} ↳ Duties – bore chn.
  & \hlblue{Heidi's Waltz} \\
Vagabonds passing
  & Tam Lin A \\
Gypsies — musicians
  & Tam Lin B \\
It was her Johnny fair / Heard he'd been married
  & \hlblue{Heidi's Waltz \textbf{highbit}} 657 \circref{D} \\
Plan together
  & \\
Happily ever after
  & \hlyellow{RTG x2 @ end} ↙ \\
\midrule
\multicolumn{2}{@{}l@{}}{\structlabel{BUT}} \\
Punished JF / Hanged {[}by{]}
  & \\
Countess's room
  & \hlpink{Hidden Valley} \\
Stairs – isolation
  & \struck{Haylands} \hspace{0.3em} \hlorange{Coconut dance (b)} \\
\bottomrule
\end{tabular}

\bigskip
\hlyellow{RTG} \textbf{to end.}


\end{document}
